% Chapter 2

\chapter{Literature Review}
\label{ch:literature_review}

This chapter reviews the research background most relevant to the thesis. The review begins with autonomous UGV navigation in unstructured terrain, then moves to UAV-assisted ground navigation, rule-based maneuvering approaches, learning-based trajectory prediction, graph-based multi-agent modeling, and communication-aware robotic coordination. The chapter concludes by identifying the specific research gap that motivates the proposed framework.

\section{UGV Navigation in Unstructured Terrain}

Autonomous UGV navigation in unstructured terrain remains a demanding research problem because perception quality, terrain geometry, and vehicle dynamics interact in ways that are difficult to model with a single sensing modality or a single control strategy. The operational difficulty of unmanned ground systems in uncertain environments is also reflected in resilience-oriented studies, where \citeauthor{krecht2023resilience} emphasize that practical UGV performance depends not only on nominal control success but also on the system's ability to remain functional under environmental and operational stress \citep{krecht2023resilience}.

Terrain understanding is a central component of this problem. \citeauthor{chen2021cnn} show that vision-proprioception fusion can improve terrain classification by combining appearance cues with vehicle-state information, thereby reducing the dependence on a single sensing source \citep{chen2021cnn}. A related direction is pursued by \citeauthor{li2025terrainfusion}, who fuse vision and vehicle dynamics for UGV terrain classification and again highlight that terrain interpretation becomes more robust when external perception is combined with motion-related signals \citep{li2025terrainfusion}. These studies are important because they move terrain reasoning beyond purely image-based classification and toward a more embodied interpretation of the driving surface.

Robustness in unstructured navigation also depends on how the vehicle represents spatial context beyond immediate terrain labels. \citeauthor{liang2025bdfusion} demonstrate that visual--LiDAR fusion can improve place recognition resilience, which is relevant because localization confidence and environmental recognition strongly affect downstream navigation performance \citep{liang2025bdfusion}. \citeauthor{carvalho2024traversability} further show that traversability analysis can be linked to mechanical effort, allowing terrain assessment and path planning to be informed by the physical burden imposed on the vehicle \citep{carvalho2024traversability}. This is especially relevant in off-road environments, where a geometrically feasible route may still be operationally undesirable if it demands excessive effort or introduces instability.

Recent work also shows increasing interest in more advanced terrain reasoning and real-time navigation support. \citeauthor{wang2024tetr} use Transformer-based traversability evaluation to assess 3D terrain for autonomous navigation, illustrating that learning-based spatial reasoning is being extended from generic perception tasks to more terrain-specific decision support \citep{wang2024tetr}. \citeauthor{han2025realtime} propose enhanced perception and memory-guided strategies for real-time UGV navigation in complex terrains, reinforcing the view that navigation quality depends not only on local sensing but also on the ability to retain useful environmental context over time \citep{han2025realtime}. At the broader survey level, \citeauthor{chen2026obstacle} show that obstacle avoidance remains a highly active topic because multi-operational environments still expose weaknesses in perception, planning, and control integration \citep{chen2026obstacle}.

Taken together, the literature on unstructured UGV navigation suggests that reliable performance requires more than local obstacle detection. It requires terrain interpretation, traversability reasoning, multi-modal sensing, and decision logic that can cope with uncertainty over both short and medium horizons. This observation is directly relevant to the present thesis, which treats navigation not only as a collision-avoidance problem but also as a short-horizon maneuver-generation problem under partial observability.

\section{UAV-Assisted Ground Navigation}

One way to reduce the sensing limitations of a ground robot is to introduce an aerial platform that can provide a broader viewpoint and complementary environmental information. \citeauthor{asadi2020integrated} demonstrate the value of integrated UGV--UAV operation for construction site data collection, showing that aerial and ground platforms can cooperate effectively when their sensing and task roles are designed to complement one another \citep{asadi2020integrated}. Although that work focuses on data collection rather than terrain-aware maneuvering, it provides an important foundation for heterogeneous robotic cooperation.

More recent work extends this cooperative perspective toward coordinated decision-making. \citeauthor{mondal2025risk} study risk-aware and energy-constrained UAV--UGV cooperative routing using attention-guided reinforcement learning, indicating that coordinated aerial-ground systems are increasingly being treated as joint decision systems rather than as loosely connected independent robots \citep{mondal2025risk}. This shift is important because it highlights the need to reason about both agent roles and interaction constraints when designing cooperative navigation frameworks.

The growing maturity of UAV-related simulation and data-generation tools also supports this research direction. \citeauthor{jarraya2025osrfodg} present an open-source ROS~2 framework for outdoor UAV dataset generation, reflecting a broader movement toward reproducible robotic experimentation and structured data collection in simulated outdoor settings \citep{jarraya2025osrfodg}. Such frameworks are particularly relevant to this thesis because they support the transition from robotic behavior design to learning-oriented dataset construction.

Despite this progress, much of the existing UAV--UGV literature still emphasizes task allocation, routing, or sensing support rather than direct maneuver guidance for the ground robot under terrain uncertainty. In many cases, the UAV contributes information at a high level, but the literature is less explicit about how that information should influence short-horizon ground motion in a continuously operating navigation loop. This gap creates a strong motivation for frameworks in which the UAV acts as a cooperative scout while the UGV remains the main maneuvering platform.

\section{Rule-Based Navigation Approaches}

Rule-based navigation approaches remain attractive because they are interpretable, modular, and comparatively easy to validate. In safety-critical or mission-oriented robotic systems, explicit maneuver rules are often preferred during early system development because they allow designers to understand exactly why the platform behaves in a certain way. \citeauthor{vincent2012combined} illustrate one useful baseline direction by combining reactive control with reinforcement learning for an autonomous tracked vehicle, showing that classical reactive behavior remains a strong starting point even when learning is later introduced \citep{vincent2012combined}.

Several strands of the literature also show how explicit control logic can be specialized for different operational needs. \citeauthor{jiang2022differential} propose a differential-steering path-tracking and torque-distribution strategy for an unmanned ground vehicle, demonstrating that vehicle-specific control structure matters for both path quality and efficiency \citep{jiang2022differential}. \citeauthor{hu2023potential} develop a potential-field-based formation-tracking approach with collision avoidance for multi-UGV systems, illustrating how rule-based interaction logic can be extended from single-vehicle navigation to coordinated multi-agent behavior \citep{hu2023potential}. \citeauthor{kim2024adaptive} further show that adaptive sensor management based on risk maps can guide how monitoring resources are allocated around a UGV, which again reflects a decision-making philosophy built from explicit operational logic \citep{kim2024adaptive}.

Rule-based approaches remain useful because they provide predictable safety behavior and clear failure interpretation. However, they also tend to rely on handcrafted priorities, thresholds, and transitions that may not generalize gracefully across varied terrain conditions. This limitation becomes especially visible when local sensing is ambiguous or when multiple factors must be balanced at once, such as obstacle proximity, terrain confidence, goal direction, and cooperative agent context. For this reason, rule-based methods are often effective as baselines or data-generation mechanisms, but they can become brittle when required to infer smoother future maneuvers under changing conditions.

This trade-off is important for the present thesis. The goal is not to discard rule-based navigation entirely, but to use it as a structured and interpretable foundation from which trajectory-learning models can later be trained. In that sense, rule-based maneuvering is treated as both an operational baseline and a teacher mechanism for subsequent data-driven learning.

\section{Learning-Based Trajectory Prediction}

Learning-based trajectory prediction reframes motion generation as a forecasting problem. Instead of relying solely on explicit handcrafted rules, the model learns to infer likely future motion from past states, environmental context, and interaction structure. Social LSTM is one of the best-known early examples of this direction, and \citeauthor{alahi2016sociallstm} show that sequential recurrent models can predict future trajectories from temporal history while accounting for interaction effects \citep{alahi2016sociallstm}. Although that work addresses human trajectory prediction, it established an influential sequence-modeling perspective that later transferred to robotic and autonomous navigation research.

For ground robotics, data-driven trajectory planning has become increasingly relevant. \citeauthor{chen2024data} propose a real-time UGV trajectory planning and control methodology using an LSTM-based architecture, demonstrating that learned sequence models can be used directly in ground-vehicle planning pipelines rather than only in offline forecasting benchmarks \citep{chen2024data}. This is particularly relevant to the present thesis because it supports the idea that a learned model can contribute to operational maneuver generation rather than merely descriptive prediction.

The rise of attention mechanisms has further expanded the scope of trajectory modeling. \citeauthor{vaswani2023attention} introduced the Transformer architecture and showed that attention can model long-range dependencies without relying entirely on recurrent computation \citep{vaswani2023attention}. \citeauthor{giuliari2020transformer} later applied Transformer networks specifically to trajectory forecasting, demonstrating that attention-based temporal modeling is also effective for future-path prediction tasks \citep{giuliari2020transformer}. These developments are important because they suggest that short-horizon maneuver prediction can benefit from richer temporal representations than those available in purely reactive systems.

Benchmark development has also shaped this field. \citeauthor{chang2019argoverse} introduced Argoverse as a large-scale tracking and forecasting benchmark with rich map context, helping establish trajectory prediction as a serious evaluation problem rather than only a model-design problem \citep{chang2019argoverse}. However, much of the benchmark-driven trajectory literature is centered on road vehicles, crowd behavior, or structured traffic scenes. This means that while trajectory prediction methods are increasingly mature, their direct transfer to cooperative UAV--UGV terrain navigation is not automatic. Off-road robotic settings involve different sensing assumptions, different uncertainty profiles, and different forms of agent interaction.

\section{Multi-Agent Learning and Graph-Based Models}

Once navigation is treated as an interaction-rich prediction problem, graph-based modeling becomes a natural candidate. \citeauthor{kipf2017semi} show that graph convolutional networks provide a principled way to learn from structured relational data, thereby moving beyond flat feature concatenation \citep{kipf2017semi}. This idea is important for robotic systems in which the meaning of a state often depends on relationships between agents rather than on the state of a single agent alone.

Temporal interaction modeling extends this idea further. \citeauthor{yu2018stgcn} propose spatio-temporal graph convolutional networks for forecasting problems, showing how graph structure and time evolution can be modeled together \citep{yu2018stgcn}. For cooperative robotic navigation, this is highly relevant because the relative geometry between a UGV and supporting UAVs changes over time, and the usefulness of that geometry depends on temporal continuity rather than isolated snapshots.

The broader trajectory-prediction literature also motivates relational modeling. \citeauthor{alahi2016sociallstm} highlighted that trajectory quality can depend on interaction context rather than only on isolated motion history \citep{alahi2016sociallstm}. In cooperative UAV--UGV settings, the same principle applies in a different form: the UGV does not move independently of the UAVs when those UAVs provide scouting support, positional context, or guidance-related information. Accordingly, graph-based representations offer a practical way to encode who the agents are, where they are relative to one another, and how those relations evolve over time.

At the same time, there is an important distinction between the graph literature and the needs of this thesis. Much graph-based trajectory work is formulated for pedestrians, traffic agents, or generic spatio-temporal forecasting settings. The present thesis instead focuses on heterogeneous aerial-ground cooperation, where the graph nodes have different physical roles and where the UGV remains the maneuvering platform while the UAVs act as supporting scouts. This makes graph modeling relevant, but it also means that the graph must be interpreted in the context of cooperative robotic functionality rather than only generic interaction forecasting.

\section{Communication-Aware Robotic Coordination}

Cooperative robotic coordination is often described as though information exchange were immediate and reliable. In practice, however, multi-robot cooperation depends strongly on the quality of the communication channel. This issue becomes especially important when an aerial robot supplies information that influences the maneuver decisions of a ground robot, because delay, jitter, and packet loss can reduce the timeliness and trustworthiness of the shared information.

The communication dimension is particularly relevant in application areas where robotic cooperation must remain dependable under adverse conditions. \citeauthor{bhatia2026integration} reinforce the importance of resilient autonomy and dependable distributed systems in broader discussions of emerging technologies in defense and related operational domains \citep{bhatia2026integration}. In such settings, evaluating navigation performance without considering communication quality can lead to overly optimistic conclusions about system robustness.

OMNeT++ provides an established simulation environment for communication-system evaluation, and \citeauthor{varga2008omnet} show why it is useful for introducing communication realism into robotic studies \citep{varga2008omnet}. Its relevance to the present thesis is not that it replaces robotics simulation, but that it complements it by making it possible to model communication degradation alongside robot behavior. This is important because UAV-assisted guidance is not only a perception problem; it is also an information-delivery problem.

Compared with the perception and control literature, communication-aware robotic coordination remains less represented in end-to-end maneuver studies. Many navigation papers evaluate sensing, planning, and control in detail, but leave communication idealized in the background. For a cooperative UAV--UGV framework, this omission is significant, because the practical usefulness of the UAV may depend not only on what it observes, but also on whether its information reaches the UGV reliably enough to affect the maneuver at the right time.

\section{Research Gap Summary}

The literature reviewed in this chapter shows that several relevant strands have advanced considerably. UGV navigation in unstructured terrain has benefited from improved terrain classification, traversability analysis, and robust sensing strategies, as reflected in the work of \citeauthor{chen2021cnn} \citep{chen2021cnn}. UAV--UGV cooperation has gained momentum as a practical and strategically meaningful form of heterogeneous autonomy, as shown by \citeauthor{asadi2020integrated} \citep{asadi2020integrated}. Learning-based trajectory prediction has matured through recurrent, attention-based, and graph-based models that can represent temporal and relational structure, as seen in the trajectory-forecasting work of \citeauthor{giuliari2020transformer} \citep{giuliari2020transformer}. Communication simulation frameworks such as OMNeT++ also provide the means to study non-ideal information exchange in a principled way, as discussed by \citeauthor{varga2008omnet} \citep{varga2008omnet}.

Nevertheless, these strands are rarely brought together into one coherent experimental workflow. Existing work does not sufficiently address the transition from a rule-based cooperative UAV--UGV maneuvering system to a learning-based short-horizon trajectory predictor trained on behavior generated by that same system. The literature also does not clearly show how such learned models behave when reinserted into live multi-robot simulation rather than evaluated only offline. In addition, communication effects are often treated separately from maneuver prediction, even though in cooperative aerial-ground systems the usefulness of shared context depends directly on network quality.

The present thesis addresses this gap by combining five elements within one framework: a rule-based UAV--UGV navigation baseline, structured simulation-based data generation, learning-based trajectory prediction, live simulation redeployment of trained model weights, and communication-aware evaluation through OMNeT++. This integrated perspective is intended to contribute not only a predictive model comparison, but also a system-level understanding of how rule-based and learning-based cooperative navigation can be related, evaluated, and compared under more realistic operational assumptions.
